\documentclass[12pt]{article}
\usepackage[a3paper, margin=1in]{geometry}
\usepackage{amsmath, amssymb, amsthm, ulem, booktabs}

\begin{document}

\section*{Ekvivalence v limitách}

\begin{enumerate}
  \item \textbf{Definice relace $\sim$}: Realné funkce $f,g$ patří do relace $\sim_c$ právě tehdy, když platí:

  \[
    \lim_{x\to c} \frac{f(x)}{g(x)} = 1
  \]

  \item \textbf{Lemma}: Bud' funkce $f: \mathbb{R} \to \mathbb{R}, g: \mathbb{R} \to \mathbb{R}$, čísla $c,A \in \mathbb{R}$
a platí následující podmínky:

\begin{itemize}
  \item $\lim_{x\to c}f(x) = A \quad (*)$
  \item $\lim_{x\to c} \frac{f(x)}{g(x)} = 1 \quad (**)$
\end{itemize}

Pak platí: $\quad \lim_{x\to c}g(x) = A$.

\medskip
\begin{quote}
\textbf{Důkaz:}
Předpokládejme funkce $f: \mathbb{R} \to \mathbb{R}, g: \mathbb{R} \to \mathbb{R}$ a čísla $c,A \in \mathbb{R}$
spjňující $(*),(**)$.

Volbou $\epsilon = 1$ z předpokladu $(**)$ pak plyne, že existuje nějaké prstencové okolí $P(c,\delta)$ takové, že
$\frac{f(x)}{g(x)}$ je definované a platí:

\[
\forall x \in P(c,\delta) : 1-1=0 < \frac{f(x)}{g(x)} < 2=1+1
\]
\[
\Longrightarrow \ \frac{f(x)}{g(x)} \neq 0 \quad (i)
\]
\[
\Longrightarrow \ f(x) \neq 0, g(x) \neq 0 \quad (ii)
\]

Pak podle věty o podílové aritmetice limit platí:

\[
\lim_{x\to c} \frac{g(x)}{f(x)} \overset{(i)}{=} \lim_{x\to c} \frac{1}{\frac{f(x)}{g(x)}} \overset{(**)}{=} \frac{1}{1} = 1
\quad (iii)
\]

Počítejme limitu funkce $g$ použitím věty o součinové aritmetice limit:

\[
\lim_{x\to c}g(x) \overset{(ii)}{=} \lim_{x\to c}\frac{g(x)f(x)}{f(x)} = \lim_{x\to c}\frac{g(x)}{f(x)} \cdot \lim_{x\to c}f(x)
\overset{(*),(iii)}{=} 1 \cdot A = A \quad\quad \square
\]
\end{quote}

\item \textbf{Důsledek lemmatu}: Bud' realné funkce $f,g$ a relace $\sim_c$, pak platí:

\[
  f \sim_c g \land \lim_{x\to c}f(x) = A \ \Longrightarrow \lim_{x\to c}g(x) = A
\]

\end{enumerate}

\end{document}
